\section{The determinant mass rate on the original polynomial source}
\label{sec:dmr-independent}

Keep the actual complete quartet polynomial
\[
 h(s)=\prod_{\lambda\in\{\rho,\bar\rho,1-\rho,1-\bar\rho\}}(s-\lambda)^m,
 \qquad \rho=\tfrac12+\delta+i\gamma,\quad
 0<\delta<\tfrac12,\quad \gamma>2,\quad m\geq1,
\]
its entire quotient $v_h=(2\xi)/h$, and its full Taylor unit.
The measures, without alteration of their masses, are
\[
 w(y)=\frac{|v_h(1/2+iy)|^2}{2\pi},\qquad m_k=w^{*k},
 \qquad d\sigma(y)=\frac{|\Gamma(1/4+iy/2)|^2}{2\pi}\,dy.
 \tag{DMR1}
\]
Write
\[
 \alpha=\pi/2,\quad p=8m-\tfrac92,\quad B=42+8m,\quad M=B/2,
 \quad M_\alpha=\int e^{\alpha y}w(y)\,dy,\quad
 \vartheta_h=\int_{-1}^1w(y)\,dy.
\]
Both displayed masses are strictly positive. The proved TW and HMR
source envelopes, with their original constants, are
\[
 \begin{split}
 &m_k(y)\leq A_k\sigma(y),\qquad
 m_k(y)\geq a_k(1+y^2)^{-M}\sigma(y),\\
 &A_k=\frac{C_h}{c_\Gamma}k^{p+1}M_\alpha^{k-1},\qquad
 a_k=\frac{c_h\vartheta_h^{k-3}e^{-\alpha(k-3)}(k-2)^{-B}}
                  {C_\Gamma 2^{B/2}},\\
 &c_\Gamma=\sqrt2 e^{-7/6},\qquad
 C_\Gamma=\sqrt{2\sqrt5}\,e^{2/3},\\
 &m_k(y)\geq c_h M_R^{k-3}e^{-\alpha|y|}
              (1+|y|+(k-3)R)^{-B}\quad
                         (|y|\geq(k-3)R),\\
 &M_R=\int_{-R}^R e^{\alpha y}w(y)\,dy,\qquad
 M_R\uparrow M_\alpha,\qquad
 0\leq M_\alpha-M_R\leq
       \frac{2C_h}{p-1}(1+R)^{1-p}.
 \end{split}\tag{DMR2}
\]
Here $\sigma(y)$ means the density of the displayed measure,
and $k\geq3$, $R\geq1$ throughout. These are the existing original
source estimates, not new hypotheses about its zeros. The calculation
below gives a consequence for every full polynomial Gram.

\subsection{A global lower measure retaining the truncated tilted mass}

Define the following finite constants and integer:
\[
 \begin{gathered}
 Y=\max(2,(k-3)R),\qquad b=\frac{c_h}{C_\Gamma 2^{3B/2}},\\
 r=\max\left(1,\left\lceil
       \frac{\log(bM_R^{k-3}/a_k)}{\log2}\right\rceil\right),\qquad
 Q(y)=(y+iY)^r(y+i)^M,\\
 d=2r+M,\qquad F(y)=y^rQ(y).
 \end{gathered}\tag{DMR3}
\]
The logarithm defining $r$ is finite. Its complete expression is
\[
 \log(bM_R^{k-3}/a_k)
 =(k-3)\{\log(M_R/\vartheta_h)+\alpha\}
                    +B\log((k-2)/2).
 \tag{DMR4}
\]
For fixed $h,R$ it follows directly that $r=O_{h,R}(k+\log k)$;
in particular $r=O_{h,R}(k)$. No uniformity in an increasing $R$
is used. The definition always gives $r\geq1$, so no zero
root block is introduced in the determinant calculation.

The literal global inequality is
\[
 \boxed{\quad
 m_k(y)\geq bM_R^{k-3}
     \left|\frac{y^r}{(y+iY)^r(y+i)^M}\right|^2\sigma(y)
       \quad(y\in\mathbb R).\quad}
 \tag{DMR5}
\]
For $|y|\geq Y$, the last lower envelope of DMR2 applies and
\[
 1+|y|+(k-3)R\leq2(1+|y|)
                           \leq2\sqrt2(1+y^2)^{1/2}.
\]
The original Gamma upper estimate gives
$e^{-\alpha|y|}\geq \sigma(y)/C_\Gamma$.
These two inequalities give
$m_k(y)\geq bM_R^{k-3}(1+y^2)^{-M}\sigma(y)$.
Multiplication of the right side by
$y^{2r}/(Y^2+y^2)^r\leq1$ proves DMR5 there.
For $|y|\leq Y$ that multiplier is at most $2^{-r}$, and DMR3
gives $bM_R^{k-3}2^{-r}\leq a_k$. The first lower estimate of DMR2
therefore proves DMR5 on the whole remaining interval, including
$y=0$. All denominators in DMR5 are nonzero on the real source line.

\subsection{The finite Gram inequality and its exact determinant}

For $N\geq0$, let $H_N^\sigma$ be the Gram of $1,y,\ldots,y^N$
in the original $\sigma$ measure, and let $H_{k,N}$ be the Gram
of that same frame in $m_k(y)\,dy$. Both forms are positive
definite, since their densities are positive almost everywhere.
Set $n=N+1$, $L=N+d$, and define the three $n$ by $n$ matrices
\[
 \begin{split}
 {\cal A}_{jl}&=\int \overline{y^{r+j}}y^{r+l}|Q(y)|^{-2}\,d\sigma(y),\\
 {\cal B}_{jl}&=\int \overline{y^{r+j}}y^{r+l}\,d\sigma(y),\\
 {\cal C}_{jl}&=\int \overline{y^{r+j}}y^{r+l}|Q(y)|^2\,d\sigma(y),
                  \qquad 0\leq j,l\leq N.
 \end{split}\tag{DMR6}
\]
The polynomial exponential tails in DMR2 make all these integrals
finite; the denominator is bounded away from zero on the real line.
The Gram matrix of the two ordered lists of functions
\[
 f_j(y)=y^{r+j}/Q(y),\qquad
 g_j(y)=y^{r+j}\overline{Q(y)}
\]
is precisely
$\left(\begin{smallmatrix}{\cal A}&{\cal B}\\
{\cal B}&{\cal C}\end{smallmatrix}\right)$.
The cross term is $ {\cal B}$ because
$\overline{f_j}g_l=\overline{y^{r+j}}y^{r+l}$.
Completing the square in its second coefficient vector proves
\[
 {\cal A}\succeq{\cal B}{\cal C}^{-1}{\cal B},\qquad
 \det{\cal A}\geq(\det{\cal B})^2/\det{\cal C}.
 \tag{DMR7}
\]
There is no interchange of the noncommuting factors in this statement.
The source inequality DMR5 now gives
\[
 \det H_{k,N}\geq (bM_R^{k-3})^n\det{\cal A}.
 \tag{DMR8}
\]

The original Gamma monic polynomials $b_j(y)$ have the exact
squared norms and recurrence proved in TW10:
\[
 \gamma_j=\sqrt{2\pi}\frac{(2j)!}{4^j},\qquad
 yb_j=b_{j+1}+j(j-\tfrac12)b_{j-1},\qquad b_0=1.
 \tag{DMR9}
\]
The absent $b_{-1}$ term is zero. Change from increasing powers to
these monic polynomials has determinant one; consequently
$\det H_N^\sigma=\prod_{j=0}^N\gamma_j$.
The exterior product of $y^r,\ldots,y^{N+r}$ has coefficient one
on $b_r\wedge\cdots\wedge b_{N+r}$. Indeed its coefficient matrix
in increasing monic degrees has diagonal entries one, and all
entries above each polynomial's degree are zero. Orthogonality
of the exterior monic frame and positivity of every squared
coefficient then imply
\[
 \det{\cal B}\geq\prod_{j=r}^{N+r}\gamma_j .
 \tag{DMR10}
\]
This is the ordinary determinant-of-Gram identity, obtained by
expanding the squared norm of that finite exterior product.

The roots of the particular auxiliary monic polynomial $F$ are
$0,-iY,-i$, with complete multiplicities $r,r,M$. Let $ {\cal J}$
be the map taking the divided Taylor derivatives
$P^{[a]}(z)=P^{(a)}(z)/a!$ at those roots in that order. Its kernel
on polynomials of degree at most $L=N+d$ is exactly
$F\mathbb C[y]_{\leq N}$. Its restriction $ {\cal V}$ to degrees
at most $d-1$ has determinant
\[
 \det{\cal V}=
  (-iY)^{r^2}(-i)^{rM}\{i(Y-1)\}^{rM},\qquad
 |\det{\cal V}|=Y^{r^2}(Y-1)^{rM}\geq1.
 \tag{DMR11}
\]
To verify the formula and all factorials, replace the derivatives
in each root block by evaluations at that root plus distinct
$\varepsilon x_a$. Divide the ordinary Vandermonde determinant
by the product of the three within-block Vandermonde factors
and by their powers of $\varepsilon$. Taylor expansion and
multilinearity give the divided derivative determinant in the
limit. The cross-block factors give exactly DMR11. In particular
$ {\cal V}$ is invertible, and $ {\cal J}$ is onto.

In the orthonormal frame $e_j=b_j/\sqrt{\gamma_j}$ the positive
Gram of the observation is
\[
 ({\cal K}_L)_{(z,a),(z',a')}=
       \sum_{j=0}^L e_j^{[a]}(z)\,
                       \overline{e_j^{[a']}(z')}.
 \tag{DMR12}
\]
This is a $d$ by $d$ matrix and includes all $r,r,M$ derivatives.
The exact determinant identity is
\[
 \boxed{\quad
 \det{\cal C}=
       \left(\prod_{j=0}^L\gamma_j\right)
          \frac{\det{\cal K}_L}{|\det{\cal V}|^2}.
 \quad}\tag{DMR13}
\]
Here is its full source proof. The ordered frame
$1,\ldots,y^{d-1},F,Fy,\ldots,Fy^N$ has coefficient determinant
one in all degrees $0,\ldots,L$. Its last block has Gram
$ {\cal C}$. In any positive source Gram $H$, the minimum-norm
vector with observation $Jv=z$ is
$H^{-1}J^*(JH^{-1}J^*)^{-1}z$; subtracting it leaves a vector
in $\ker J$ orthogonal to the displayed minimum vector.
Its squared norm is $z^*(JH^{-1}J^*)^{-1}z$.
Apply this identity to $ {\cal J}$. In the first $d$ remainder
coordinates, the Schur-complement Gram is
$ {\cal V}^*{\cal K}_L^{-1}{\cal V}$.
Taking the determinant of the block Gram therefore gives
$\prod_{j=0}^L\gamma_j=\det{\cal C}\,
|\det{\cal V}|^2/\det{\cal K}_L$, which is DMR13.
Thus the auxiliary kernel is used through an exact map on the
original Gamma source; it is not an asserted replacement of
the original counterfactual packet.

\subsection{A uniform finite bound on every auxiliary divided jet}

The recurrence in DMR9 implies the convergent generating function
\[
 \sum_{j=0}^\infty \frac{b_j(z)}{j!}t^j
       =(1+t^2)^{-1/4}\exp(z\arctan t),\qquad |t|<1.
 \tag{DMR14}
\]
Indeed multiplication of the recurrence by $t^j/j!$ and summation
gives $(1+t^2)\partial_t G=(z-t/2)G$, $G(0,z)=1$.
Solving this analytic equation yields the right side; equality
of its Taylor coefficients proves the asserted identity.
All branches are the analytic branches with value one, or zero
for $\arctan$, at $t=0$.

Set
\[
 Z_{L,Y}=\frac{e^2(L+1)(L+2)(2L+1)}{\sqrt{2\pi}}\,
                  [2(L+2)]^{Y+1}.
 \tag{DMR15}
\]
Then every diagonal entry in DMR12 is at most $Z_{L,Y}$, and
\[
 \det{\cal K}_L\leq Z_{L,Y}^{\,d}.
 \tag{DMR16}
\]
For a complete coefficient proof, put $s=(L+1)/(L+2)$.
On $|t|=s$ the absolutely convergent arctangent series gives
\[
 |\arctan t|\leq\operatorname{arctanh}s
 =\tfrac12\log((1+s)/(1-s))
 \leq\tfrac12\log(2(L+2)).
\]
Also $|1+t^2|\geq(1-s)^2$ and $s^{-j}\leq e$ for $j\leq L$:
the latter follows from
$j\log(1+1/(L+1))\leq j/(L+1)\leq1$.
Cauchy's coefficient formula in $t$ yields, for $|z|\leq Y+1$,
\[
 \frac{|b_j(z)|}{j!}
 \leq e(L+2)^{1/2}[2(L+2)]^{(Y+1)/2}.
 \tag{DMR17}
\]
Apply Cauchy's derivative formula on the circle of radius one
about any of $0,-iY,-i$. These circles lie in $|z|\leq Y+1$;
for every divided derivative order $a$ the same right side
bounds $|b_j^{[a]}(z_0)|/j!$. No factorial is dropped, because
the derivatives are divided by $a!$ and the circle radius is one.
Finally
\[
 \frac{(j!)^2}{\gamma_j}
 =\frac{4^j}{\sqrt{2\pi}\binom{2j}{j}}
 \leq\frac{2j+1}{\sqrt{2\pi}} .
 \tag{DMR18}
\]
The inequality follows by summing the $2j+1$ binomial coefficients,
each of which is at most the central one. Equations DMR17--18
bound each squared divided jet by $Z_{L,Y}/(L+1)$.
Summation in DMR12 proves the diagonal bound. The determinant
bound follows by orthogonally projecting each successive row
of its Gram frame away from the earlier rows: each resulting
squared length is no greater than the original diagonal
entry, and their product is its determinant.

\subsection{All finite factorial losses and the determinant limit}

Put $g_0=\log\sqrt{2\pi}$ and define the nonnegative quantity
\[
 \begin{split}
 {\cal E}_{N,r,Y}={}&2r^2\log(N+2r)
     +M\{g_0+2L\log L\}\\
     &+2rg_0+4r^2\log(\max(1,r))
       +(2r+M)\log Z_{L,Y},
             \qquad L=N+2r+M.
 \end{split}\tag{DMR19}
\]
Since $M\geq25$, $r\geq1$ and $Y\geq2$, all its logarithms
used as upper bounds are positive. The following two finite
inequalities hold for every $N\geq0$:
\[
 \boxed{\quad
 n\{\log b+(k-3)\log M_R\}-{\cal E}_{N,r,Y}
 \leq
       \log\frac{\det H_{k,N}}{\det H_N^\sigma}
 \leq n\log A_k,\qquad n=N+1.
 \quad}\tag{DMR20}
\]
For the upper side apply the first inequality of DMR2 to every
coefficient vector and take determinants.
For the lower side combine DMR7--8, DMR10, DMR13,
DMR11's $|\det{\cal V}|\geq1$, and DMR16. The remaining
logarithm of Gamma norm products is precisely
\[
 T=2\sum_{j=r}^{N+r}\log\gamma_j
          -\sum_{j=0}^{N+2r+M}\log\gamma_j
          -\sum_{j=0}^N\log\gamma_j .
 \tag{DMR21}
\]
Writing each sum as a difference of two prefix sums proves,
without any assumption about whether $r\leq N$, the identity
\[
 T=\sum_{j=1}^r
      \{\log\gamma_{N+j}-\log\gamma_{N+r+j}\}
       -\sum_{j=N+2r+1}^{N+2r+M}\log\gamma_j
       -2\sum_{j=0}^{r-1}\log\gamma_j .
 \tag{DMR22}
\]
The exact ratio
$\gamma_{j+1}/\gamma_j=(j+1)(j+1/2)$ gives
\[
 \log\gamma_{N+r+j}-\log\gamma_{N+j}
 \leq2r\log(N+2r)\qquad(1\leq j\leq r).
\]
The factorial formula in DMR9 gives, for every $j\geq0$,
\[
 \log\gamma_j\leq g_0+2j\log(\max(1,j));
\]
for $j\geq1$ this follows from $(2j)!\leq(2j)^{2j}$ and
$j\log4=2j\log2$, and for $j=0$ it is equality.
Apply this last bound to the $M$ terms at the top and to
the $r$ terms at the bottom of DMR22. It follows that
$T-d\log Z_{L,Y}\geq-{\cal E}_{N,r,Y}$, proving DMR20
with every finite factorial loss accounted for.

Now fix the original packet $h$ and take exactly the original
degrees
\[
 q=(k+1)^2[1+k(m-1)],\qquad
 N\in\{q-1,q,2q-1,2q\}.
 \tag{DMR23}
\]
For fixed $R$, DMR4 gives $r=O_{h,R}(k)$, DMR3 gives
$Y=O_R(k)$, and $L=N+O_{h,R}(k)$. DMR15 gives
$\log Z_{L,Y}=O_R(k\log(N+k+2))$. Thus the fully displayed
expression DMR19 satisfies
\[
 {\cal E}_{N,r,Y}
     =O_{h,R}\bigl(N\log(N+k+2)+k^2\log(N+k+2)\bigr).
 \tag{DMR24}
\]
Here $N+1\geq q\geq(k+1)^2$ and $N=O_h(k^3)$, so
$ {\cal E}_{N,r,Y}/(k(N+1))\to0$ at each of the four degrees.
Moreover
\[
 \frac{\log A_k}{k}
  =\frac{k-1}{k}\log M_\alpha
      +\frac{(p+1)\log k+\log(C_h/c_\Gamma)}{k}
       \longrightarrow\log M_\alpha .
 \tag{DMR25}
\]
Divide DMR20 by $k(N+1)$, first keep $R$ fixed and let
$k\to\infty$. The lower limit is at least $\log M_R$ and
the upper limit is at most $\log M_\alpha$. Then let
$R\to\infty$ in the former statement; DMR2 supplies
$M_R\uparrow M_\alpha>0$. This proves
\[
 \boxed{\quad
 \lim_{k\to\infty}
  \frac{1}{k(N+1)}
       \log\frac{\det H_{k,N}}{\det H_N^\sigma}
       =\log M_\alpha
   \quad\text{at all four exact degrees in DMR23}.
 \quad}\tag{DMR26}
\]
The order of the two limits is explicit. No pointwise lower
bound for a one-factor zeta numerator has been introduced.

Return now to the original coordinate $S=k/2+iy$.
The coefficient map $P(S)\mapsto P(k/2+iy)$ is invertible
on every degree-$N$ source and has determinant
$\prod_{j=0}^Ni^j$ of absolute value one. It acts on both
of the original Gram forms by the same congruence and
leaves their determinant ratio exactly unchanged.
Consequently DMR20 and DMR26 hold literally for the
original $S$-coordinate Grams. The cyclic observation
and its full Taylor unit are retained by applying their
same original coefficient matrices on this source;
neither was used as a surrogate for the polynomial
matrix whose determinant has just been computed.

\subsection{The positive total spectral deficit}

For the original $S$-coordinate matrices put
\[
 D_{k,N}=(H_N^{\sigma,k})^{-1/2}H_{k,N}
                              (H_N^{\sigma,k})^{-1/2},\qquad
 \Xi_{k,N}=(N+1)\log A_k-\log\det D_{k,N}.
 \tag{DMR27}
\]
This is a congruence with the explicitly displayed original
Gamma metric, not a change in either original measure.
DMR2 proves $0<D_{k,N}\preceq A_k I$ and hence
$\Xi_{k,N}\geq0$. Equations DMR25--26 give
\[
 \boxed{\quad \Xi_{k,N}=o(kq)\quad
       (N=q-1,q,2q-1,2q).\quad}
 \tag{DMR28}
\]
The exact finite upper bound supplied by DMR20 is
\[
 0\leq\Xi_{k,N}
  \leq(N+1)\{\log A_k-\log b-(k-3)\log M_R\}
                          +{\cal E}_{N,r,Y}.
 \tag{DMR29}
\]
The full quotient and retained-kernel propagation of this
nonnegative total deficit is established by the accompanying
explicit source/quotient Schur-complement calculation.
It does not propagate a condition-number lower bound to
a quotient, and it does not assign a sign to a difference
of four endpoint logarithms.

\subsection{A single explicit tensor-order rate}

The finite bound DMR29 permits an increasing truncation radius.
All constants needed for this step can be kept explicit.
For $R\geq1$ positivity gives $M_1\leq M_R\leq M_\alpha$.
Evenness of $w$ also gives $M_\alpha\geq\vartheta_h$.
Define the following constants using only the original packet
and the original envelope constants:
\[
 \begin{gathered}
 C_r=1+\frac{\alpha+\log(M_\alpha/\vartheta_h)+B}{\log2},
 \qquad C_d=2C_r+M/3,\qquad C_L=1+C_d,\\
 C_E=6C_r^2+(2M+4)C_d/3,\qquad
 C_s=\left|\log\frac{C_h}{c_\Gamma b}+2\log M_\alpha\right|,\\
 T_k=\log(4m(C_L+2))+3\log(k+2),\qquad
 C_{\rm tail}=\frac{2C_h}{(p-1)M_1}.
 \end{gathered}\tag{DMR30}
\]
In particular none of these constants depends on $R$.
The elementary bound $\lceil x\rceil\leq x+1$ and DMR4 give
\[
 1\leq r\leq C_r k,\qquad
 d\leq C_d k,\qquad Y\leq kR,\qquad L\leq C_L N
 \quad(N\in\{q-1,q,2q-1,2q\}).
 \tag{DMR31}
\]
For the first inequality use
\[
 \log(bM_R^{k-3}/a_k)
 \leq k\{\alpha+\log(M_\alpha/\vartheta_h)\}+B\log k,
 \qquad \log k\leq k.
\]
The right side is positive, so it also bounds the maximum with
one in DMR3 after the stated ceiling estimate. For the second
use $M\leq Mk/3$; for the third use $2\leq kR$.
For the last use $N\geq q-1\geq k^2$, so
$L=N+d\leq N+C_dk\leq(1+C_d)N$.
These arguments include $k=3$.
Also $N\leq2m(k+1)^3$, which implies
\[
 \log(2(L+2))\leq T_k,\qquad
 \log(N+2r)\leq T_k,\qquad \log r\leq T_k .
\]
From the exact formula DMR15,
$\log Z_{L,Y}\leq2+(Y+4)T_k$:
its three polynomial logarithms are each at most $T_k$,
its exponential term contributes at most $(Y+1)T_k$,
and $-\log\sqrt{2\pi}<0$.
Substituting these estimates in the exact DMR19 proves
\[
 {\cal E}_{N,r,Y}
 \leq 2MNT_k+k^2(C_E+C_dR)T_k+(g_0+2)C_dk .
 \tag{DMR32}
\]
Indeed the two quadratic $r$ terms contribute at most
$6C_r^2k^2T_k$. The $2MLT_k$ term is at most
$2MNT_k+2MC_dkT_k$. The remaining $d(Y+4)T_k$ term
is at most $C_dk^2RT_k+4C_dkT_k$.
Using $k\leq k^2/3$ gives exactly $C_E$ for those two
linear terms. Finally the remaining $g_0d+2d$ term
is at most $(g_0+2)C_dk$.

The explicit finite bound at all four original endpoints is
\[
 \boxed{\begin{split}
 0\leq\frac{\Xi_{k,N}}{kq}
 \leq\left(2+\frac1q\right)\bigg\{&
 C_{\rm tail}(1+R)^{1-p}
 +\frac{(p+1)\log k+C_s}{k}\\
 &+\frac{(2M+C_E+C_dR)T_k}{k}
 +\frac{(g_0+2)C_d}{(k+1)^2}\bigg\}.
 \end{split}}\tag{DMR33}
\]
To prove every scalar term, the exact expression in DMR29 is
\[
 \begin{split}
 \log A_k-\log b-(k-3)\log M_R
  ={}&k\log(M_\alpha/M_R)+(p+1)\log k\\
     &+\log(C_h/(c_\Gamma b))-\log M_\alpha+3\log M_R.
 \end{split}
\]
Since $M_R\leq M_\alpha$, its final three terms are at most
$C_s$. This is valid whether the masses are below or above one.
Also
\[
 0\leq\log(M_\alpha/M_R)
 \leq(M_\alpha-M_R)/M_R
 \leq C_{\rm tail}(1+R)^{1-p}.
\]
Here $\log(1+x)\leq x$ follows by integrating
$(1+t)^{-1}\leq1$ on $0\leq t\leq x$.
Divide DMR32 by $k(N+1)$; the facts
$k^2/(N+1)\leq1$ and $N+1\geq(k+1)^2$ give the other
three terms in braces. Finally $(N+1)/q\leq2+1/q$,
proving DMR33.

For $k\geq3$ choose the particular radius
\[
 R_k=\left(\frac{k}{\log(k+2)}\right)^{1/p}>1,\qquad
 \eta=\frac{p-1}{p}>0.
 \tag{DMR34}
\]
The inequality $k>\log(k+2)$ follows at $k=3$ and persists
because the derivative of $x-\log(x+2)$ is positive.
Thus every earlier finite inequality applies to this radius.
Writing $x_k=\log(k+2)/k\in(0,1)$, one has
$(1+R_k)^{1-p}\leq R_k^{1-p}=x_k^\eta$ and
$R_k x_k=x_k^\eta$. Also $T_k=O_h(\log(k+2))$ and
each remaining term in DMR33 is $O_h(x_k)$.
Since $x_k\leq x_k^\eta$, DMR33 gives the single quantified rate
\[
 \boxed{\quad
 0\leq\frac{\Xi_{k,N}}{kq}
   =O_h\left(\left(\frac{\log(k+2)}{k}\right)^{(p-1)/p}\right),
 \quad N=q-1,q,2q-1,2q.
 \quad}\tag{DMR35}
\]
DMR33, with the specified $R_k$ and constants DMR30,
is the complete finite bound behind this notation.

\subsection{Spectral counting and the common original ambient source}

Let $\lambda_1,\ldots,\lambda_{N+1}$ be the eigenvalues
of the original relative Gram operator $D_{k,N}$, with their
algebraic multiplicities. Positivity makes them real and positive,
and DMR2 bounds them by $A_k$. For every $\epsilon>0$,
\[
 \boxed{\quad
 \#\{j:\lambda_j/A_k\leq e^{-\epsilon k}\}
       \leq\frac{\Xi_{k,N}}{\epsilon k}.
 \quad}\tag{DMR36}
\]
Indeed $\Xi_{k,N}=\sum_j-\log(\lambda_j/A_k)$ is a sum
of nonnegative terms, and every counted term is at least
$\epsilon k$. Thus DMR35 also bounds the fraction of such
eigenvalues by a constant times
$\epsilon^{-1}(\log(k+2)/k)^{(p-1)/p}$ at the actual
four source dimensions. This statement retains the full
spectrum and does not pass a condition-number lower bound
through an unproved quotient map.

For completeness, the one common original ambient source is
$\mathcal P_{2q}$, with both original forms. At any of the
four endpoints let
$C_N:\mathcal P_N\hookrightarrow\mathcal P_{2q}$
be the increasing-coefficient inclusion. Its entries are
$(C_N)_{ab}=1$ for $a=b\leq N$ and zero otherwise.
Then literally
$H_{k,N}=C_N^*H_{k,2q}C_N$ and
$H_N^{\sigma,k}=C_N^*H_{2q}^{\sigma,k}C_N$.
The exact restriction determinant identity QDC12--16 gives
$0\leq\Xi_{k,N}\leq\Xi_{k,2q}$.
One may verify its use directly by taking
$V=(H_{2q}^{\sigma,k})^{1/2}C_N
       (H_N^{\sigma,k})^{-1/2}$, for which $V^*V=I$.
Compression of $D_{k,2q}/A_k$ to this isometry has determinant
$\det D_{k,N}/A_k^{N+1}$; its complementary Schur determinant
lies in $(0,1]$. Taking negative logarithms gives the stated
inequality with the original source dimensions retained.

Let $\Delta F^\sigma_k=F_{\rm ar,k}-F_{\sigma,k}$ be the
original cyclic quotient four-endpoint sum defined in QDC23--24,
and let $\Delta F^K_k,\Delta F^B_k$ be its retained-kernel
and actual-boundary sums in QDC28--32. Those exact maps give
\[
 \boxed{\quad
 |\Delta F^\sigma_k|,\ |\Delta F^K_k|,\ |\Delta F^B_k|
      \leq2\Xi_{k,2q}.
 \quad}\tag{DMR37}
\]
In fact each of their four nonnegative quotient or restriction
deficits is at most its endpoint $\Xi_{k,N}$, and the exact
signs are $(-,-,+,+)$ on those deficits. Each positive side
and each negative side is therefore at most $2\Xi_{k,2q}$.
The common scalar $q\log A_k$, or the corresponding actual
kernel or boundary dimension times $\log A_k$, cancels with
the four signs; no scalar attached to a different source
dimension is canceled. Equations DMR35 and DMR37 prove the
same quantitative $o(kq)$ control for all three actual sums
relative to the stated fixed one-factor Gamma reference.
The exact comparison with the separate original $k$-fold
Gamma reference is a further retained calculation. No sign
or RH endpoint is inferred from DMR37.
